\chapter{The positive-unit $e$ subdomain of weak-C2 R4}
\label{component:p78dfeedee062}
\noindent\textbf{Component }\path{p78dfeedee062}\par
{\small\noindent Source: \path{R17/proofs/c1_source_ends/r4_e_unit_all_tests.tex}\par}
{\small\noindent SHA256: \path{78dfeedee06220f42324bb95b6cd5822ad05cf45d1a52af808765ccc0cd70f14}\par}
{\small\noindent Inventory use: R4E1-R4E16\par}
\medskip
\section{The positive-unit $e$ subdomain of weak-C2 R4}
\label{p78dfeedee062:sec:r17-r4-eunit}

\begin{theorem}
Keep the original physical stationary C2 source, $d=1$, literal
$\lambda_{\rm event}=\beta=0$, $n\ge0$, all strict source gaps, the
positive-infinitesimal total-scale intersection, and every inherited
prior-admission exclusion. Suppose the ordered source is in R4 and
$e=1-r$ is a positive bounded unit. Then the entire original four-state
parent is a finite tensor SOS for every independent projective testing
pair and every real or complex state. This admits precisely the
$e$-unit subdomain of R4; the $e$-infinitesimal part is not claimed.
\end{theorem}

\paragraph{Source entry.}
Use the exact WE1 source parameters
\[
 a=z/B,\quad t=y/B,\quad h=H/B,\quad l=1-h,
 \quad g=1-ae+haB,\quad j=1+aBr-l/t.
\]
On R4, $B,z,y,j$ are positive infinitesimals and $t,h,l$ positive
bounded units. Since $e$ is a unit, $g>0$ forces $a$ bounded.
The exact $j$ identity gives $t_0=l_0$. In the following residue
formulas the letters on the right denote their real residues, with
$h+l=1$, $e=1-r$, and $g=1-ae\ge0$. Put
$M_*=1+ar-ael=g+a(r+eh)>0$. Joint regular reduction of WE1 gives
\[
 (nr)_0={e\over2(1+a)},\quad
 (mg)_0={a^2e\over2(1+a)},\quad
 (Ah)_0={l(1+ar)\over2(1+a)},\quad
 (Dj)_0={1+ar\over2l(1+a)},                         \tag{R4E1}
\]
\[
 C_0={M_*\over2l(1+a)},\qquad
 \alpha_{*,0}:=(\alpha/B^2)_0={(1+ar)M_*\over2(1+a)}.\tag{R4E2}
\]
Thus $A,C,\alpha_*,Dj$ are positive bounded units, while $D$ is
unbounded. The coefficient $n$ is either a positive unit or unbounded,
according as $r$ is a unit or infinitesimal. In particular literal
$n=0$ is impossible here by R4E1. No coefficient boundary was discarded.
If $g$ is a unit then $\mu=m/a^2$ is a positive bounded unit with
\[
                 \mu_0={e\over2g(1+a)}.              \tag{R4E3}
\]
If $g$ is infinitesimal, $a$ is a unit, $m$ unbounded, and $mg$ a
positive unit. These are joint regular formulas: after clearing $r,g,j$
by the displayed products, the remaining WE1 denominators are products
of $t,l,1-hB,\Delta$, with $\Delta_0=1+a>0$. No relative gap rates
are imposed. Write $N_*=n+1/2$, so $N_*^{-1}$ and $N_*^{-1/2}$ are
bounded and $nr$ is bounded in all source cases.

\paragraph{Exact removal of the divergent D channel.}
Use the complete WA3--WA8 identity in the scaled roots
\[
 {\cal R}=(\xi,\eta,R_1,R_2',T)^T,
 \quad \xi=X+ruZ,\quad\eta=Y-huW,
\]
\[
 R_1=uW-vZ/t,\quad R_2'=avW-uZ,\quad T=uvW-BZ,
 \qquad F/4={\cal R}^TG{\cal R}.
\]
The exact full-matrix identity RD3--RD6, proved independently of any
small-$h$ hypothesis in the R17 small-$r$ R8 lemma, gives
\[
 L_D=uv\xi-\eta+lR_1+ruvR_2'-arvT
     =uW+uvX-Y-vZ+jvZ,
\]
\[
 F/4=D L_D^2+{\cal R}^T(G|_{D=0}+DjH_D){\cal R}.       \tag{R4E4}
\]
Here the only nonzero symmetric entries of $H_D$ are
$(13)=tuv$, $(15)=-tu$, $(34)=rtuv$, $(35)=-artv$.
The underlying root syzygy is exact, and this identity contains the
full hard block and both complete first normals. Formal $D=0$ does
not assert a physical modified source. Put $\widehat D=\sqrt D L_D$
and solve the source-unit coefficient $l$:
\[
 R_1={1\over l}\left(D^{-1/2}\widehat D+\eta
       -uvN_*^{-1/2}\widehat\xi-ruvR_2'+arvT\right),
 \qquad\widehat\xi=\sqrt{N_*}\xi.                  \tag{R4E5}
\]
Every coefficient is bounded and separated, with no mixed inverse.

\paragraph{All local tests when g is a unit.}
The $D$-free matrix after the $\widehat\xi$ normalization has bounded
pure-source coefficients. Indeed all scalar occurrences of $n$ are
through $nr^2$, and all normal occurrences through $nr$; the hard
coefficient $n/N_*$ is bounded. $A,C,m,\alpha_*,\mu$ are bounded,
and $t,l,g$ have unit inverse. The only possible $a$ pole after using
$R_2'$ is
\[
 Q_*={b_2-\alpha_*\over a},\quad
 b_2=2a^2\mu h+Cl^2+(Ah)h+(B^2\alpha_*+1/2)h^2.     \tag{R4E6}
\]
The exact WE3 numerator vanishes at $a=0$ for arbitrary $B,h,r,t$.
After cancellation its denominator has only unit factors on this
$g$-unit chart; no $D$ or $j$ occurs. All other scaled WA coefficients
have the same regular rewrites as RD7. Thus R4E5 and bounded $Dj$
give a bounded full five-root Gram on
$(\widehat\xi,\eta,\widehat D,R_2',T)$.

At the local testing residue $u=v=0$, $D^{-1/2}$ vanishes and
$R_1=\eta/l$. The complete residue has independent pivots
$A_0/l_0^2$, $1$, $\alpha_{*,0}$ on $\eta,\widehat D,T$.
Its remaining $\widehat\xi,R_2'$ block is
\[
 \begin{pmatrix}
 1+a_0^2\mu_0\nu&-a_0\mu_0\sqrt\nu\\
 -a_0\mu_0\sqrt\nu&\mu_0
 \end{pmatrix},\qquad
 \nu=(N_*^{-1})_0\ge0.                              \tag{R4E7}
\]
The determinant is exactly $\mu_0>0$. This includes $r$ infinitesimal
($\nu=0$), $a$ infinitesimal, and every independent relative rate.
For reference, the $\eta$ pivot is the exact cancellation
\[
 (A+C+1/2)+{2[-(C+1/2)l+Ah]\over l}
       +{(C+1/2)l^2+Ah^2\over l^2}=A/l^2.
\]
No scalar or normal term is omitted in this computation.

\paragraph{All local tests when g is infinitesimal.}
Here $a$ is a unit. Use the unscaled $R_2=vW-uZ/a$, substituting
$R_2'=aR_2$ in R4E4--R4E5, and the exact
centered-$m$ matrix identity R5.2--R5.4 before R4E4. Namely, retain
\[
 (m-K)R_{mc}^2,\qquad R_{mc}=-\xi+uv\eta+(1-g)R_2+huT,
\]
and replace the remaining matrix by
$G|_{m=K}+\tau H_m$, where $\tau=(m-K)g$ is bounded with positive
residue and
\[
 (H_m)_{14}=-1,\quad(H_m)_{25}=v,\quad
 (H_m)_{44}=2-g,\quad(H_m)_{45}=hu.
\]
Choose a sufficiently large positive real $K$; $m-K>0$ for every
such $K$ since $m$ is unbounded. This identity and R4E4 commute,
as $H_m$ has no $D$ dependence. All remaining coefficients are
bounded: $a$ is a unit, $m$ has been replaced by $K$, $\tau$ is
bounded, and the only unbounded $n,D$ are treated by R4E5.
The same three residue pivots remain. The other block is now
\[
 \begin{pmatrix}
 1+K\nu&-(K+\tau_0)\sqrt\nu\\
 -(K+\tau_0)\sqrt\nu&K+2\tau_0
 \end{pmatrix},\qquad
 \det=K+2\tau_0-\tau_0^2\nu.                         \tag{R4E8}
\]
Since $0\le\nu\le2$, any real $K>2\tau_0^2+1$ makes it strict.
No source opening is compared with a testing rate. Restore the exact
retained square. Equations R4E7--R4E8 and WA10 therefore prove the
entire local parent for all independent infinitesimal or zero $u,v$,
with arbitrary states and either sign. The finite separated lift
whitens the strict pure-source matrix and handles its finitely many
infinitesimal testing monomials; it introduces no mixed radicals.

\paragraph{Finite remote tests.}
If both finite angles have nonzero residue, use the inherited
all-noncommuting-residues theorem, at arbitrary coefficient scales
(the precise entry is the inherited \path{penpaper_noncommuting.tex}).
Its hypotheses apply to this unchanged physical parent. It remains
to give the two one-unit-angle charts. In these remote charts,
positive original event coefficients may be clipped at positive real
lower bounds and their exact positive excess packets retained; no
contact assertion is made for the clipped polynomial.

If $u$ is a bounded unit and $v$ infinitesimal or zero, clip $n,A,C$
at positive real lower bounds $n_*,A_*,C_*$, and retain all their
excesses, the entire $m$ and $\alpha$ packets. Normalize the original
$D$ event by $\widehat J=\sqrt D(uW+uvX-Y-vZ)$ and solve
\[
 W={D^{-1/2}\widehat J+Y+vZ-uvX\over u}.
\]
The whole remainder has bounded coefficients on
$(\widehat J,X,Y,Z)$. Its residue is
\[
 \widehat J^2+n_*D_uX^2+\tfrac12(X+uZ)^2+A_*D_uY^2\succ0,
 \qquad D_u=1+u_0^2.                                \tag{R4E9}
\]
Here and in the remaining residue displays the tests denote their real
residues. Positivity is strict because $u_0\ne0$. All source cases,
including unbounded $n,m$, have been covered by exact excess packets.

If $u$ is infinitesimal or zero, $v$ a bounded unit, and $a$ a unit,
$m$ has a positive real lower bound even when it is unbounded.
Clip $n,m,A,C$ at positive real bounds and retain their excesses and
$\alpha$. Solve instead
$Z=(uW+uvX-Y-D^{-1/2}\widehat J)/v$. The bounded remainder residue is
\[
 \widehat J^2+(n_*+1/2)D_vX^2+m_*(vW-X)^2
                  +(A_*+C_*+1/2)D_vY^2\succ0.       \tag{R4E10}
\]
This is strict since $m_*>0$ and $v_0\ne0$.

Finally let $a$ be infinitesimal in this chart. Then $g_0=1$,
$\mu=m/a^2$ is a unit, and
\[
       A_0={l\over2h},\qquad C_0={1\over2l}.          \tag{R4E11}
\]
Keep $S=aW$ as an extra legal root and write $R=uW$ everywhere
else. Clip only $n$ at a positive real $n_*$, retaining its excess
and the $\alpha$ packet. Retain the actual $A,C,m,D$ and the whole
signed core. Use
\[
 R=Y+vZ-uvX+D^{-1/2}\widehat J,\qquad
 mM^2=\mu(vS-aX+auvY-auZ)^2.                         \tag{R4E12}
\]
There is no inverse of $a$ or $u$ here. The complete remainder is
a bounded matrix on $(S,X,Y,\widehat J,Z)$. Its separate residue
pivots are $\mu_0V$, $(n_*+1/2)D_v$, and $1$ on
$S,X,\widehat J$. The remaining $Y,Z$ matrix is
\[
 \begin{pmatrix}
 (l+V)/(2h)&-v/2\\
 -v/2&V(h+V)/(2l)
 \end{pmatrix},\qquad
 \det={V^2(1+V)\over4hl}>0,
 \quad V=v_0^2>0.                                   \tag{R4E13}
\]
This exact computation retains the entire signed core; it is not a
state-slice argument. The possibly dependent extra root $S$ causes
no difficulty for a positive Gram. WA10 at the real angle residues
now proves each complete finite remote chart.

\paragraph{Projective endpoints and the join.}
The source bounds give positive real lower bounds for $n,C,A,D$.
Clip them at $n_*,C_*,A_*,D_*>0$, retain their whole positive
excess packets, and retain the entire original $m,\alpha$ packets.
All remaining source coefficients are real constants. The following
three exact homogeneous charts give the complete endpoint join,
without any further source hypothesis.

For $u$ unbounded put $p=1/u$ and use $p^2F/4$. If $v$ is bounded,
the selected residue at $p=0$ is
\[
 n_*D_vX^2+C_*D_vW^2+A_*D_vY^2
       +(D_*+1/2)(W+vX)^2+\tfrac12(vY-Z)^2\succ0.     \tag{R4E14}
\]
It is strict on $(X,W,Y,Z)$ for every real $v$, including zero.
The full homogeneous selected matrix has bounded polynomial entries.

For $v$ unbounded put $q=1/v$ and use $q^2F/4$. If $u$ is bounded,
write $R=uW$; every remaining occurrence of $W$ is through $R$.
The selected residue at $q=0$ is
\[
 (n_*+1/2)D_uX^2+C_*(R-Y)^2
       +(A_*+1/2)D_uY^2+D_*(uX-Z)^2\succ0.           \tag{R4E15}
\]
It is strict on $(X,R,Y,Z)$ for every real $u$, including zero,
and the exact selected polynomial uses no inverse of $u$.

If both tests are unbounded use $p=1/u,q=1/v$ and $p^2q^2F/4$,
with the legal roots $P=pZ,Q=qZ$. The exact half-core becomes
\[
 \tfrac12\{(qW+X-pqY-P)^2+(1+p^2)(Y-Q)^2
                       +(1+q^2)(p^2X^2+2XP-P^2)\}.
\]
Together with the selected event packets its residue is
\[
 (n_*+1/2)X^2+C_*W^2+A_*Y^2+D_*(X-P)^2
                                  +\tfrac12(Y-Q)^2\succ0,       \tag{R4E16}
\]
strict on $(X,W,Y,P,Q)$. Dependent roots do not affect positivity.
All finite reciprocal deviations are pure testing infinitesimals,
so the separated strict-Gram lift applies in all three charts.
Restore the original positive excesses and $m,\alpha$ packets.
Homogeneous restoration uses pure-testing squares off the endpoints;
literal $p=0$ or $q=0$ uses the polynomial endpoint value itself.

The finite local, finite noncommuting, both finite one-unit-angle
charts, and the homogeneous infinity charts exhaust all independent
tests. All four original state coordinates and both normal directions
were retained. Polynomial axis values, sign gauges, and application to
real and imaginary state parts supply the literal axes, both signs,
and complex states. The result is an all-tests admission of the
positive-unit-$e$ source subdomain of R4, with its original physicality,
strict gaps, coefficient signs, and residual intersections unchanged.
